\documentclass[10pt]{article}
\usepackage[margin=1in]{geometry}
\usepackage{amsmath,amssymb,booktabs,graphicx,microtype,hyperref}
\usepackage{siunitx}
\usepackage{xcolor}
\graphicspath{{figures/}}
\title{ORBIT: Function-Space Optimization for Rotary Query--Key Interactions}
\author{Alexandru Pop}
\date{}

\IfFileExists{generated_results.tex}{\input{generated_results.tex}}{%
  \newcommand{\OrbitConfirmLoss}{TBD}
  \newcommand{\OrbitConfirmSD}{TBD}
  \newcommand{\MuonConfirmLoss}{TBD}
}

\begin{document}
\maketitle

\begin{abstract}
Modern matrix optimizers treat query and key projection matrices as ordinary parameter
operators even though rotary attention couples them through position-dependent bilinear
interactions. We study ORBIT, an optimizer-only method that uses tiny frequency-resolved
$2\times2$ metrics to precondition query/key updates according to their predicted movement
in rotary attention-logit space. The model architecture and inference graph remain
unchanged. This manuscript is intentionally claim-free until the frozen multi-seed campaign
is complete; numerical claims are injected only from the recorded result artifacts.
\end{abstract}

\section{Introduction}
The central question is whether an optimizer for rotary attention should measure a query/key
step in raw parameter space or in the function induced on attention logits. ORBIT tests the
second hypothesis while retaining ordinary Transformer weights and standard inference.
The paper is structured to separate three questions: (i) does the functional metric improve
optimization, (ii) is the RoPE/phase structure itself necessary, and (iii) does any gain
survive held-out seeds, longer training, and increased model scale?

\paragraph{Contributions.}
The intended contributions, subject to the final prior-art and empirical gates, are:
\begin{enumerate}
  \item a cheap RoPE-aware local metric for joint query/key optimization based only on
  $2\times2$ frequency-pair statistics;
  \item an optimizer implementation that preserves the total Q/K update norm after
  functional preconditioning and changes no inference computation;
  \item a paper-grade comparison against tuned AdamW, Muon, NorMuon, AdaMuon and the
  inherited ASTRO baseline with fixed corpus, scheduling and decay policy;
  \item ablations that isolate relative-position rotations, off-diagonal phase coupling,
  and the complete functional preconditioner.
\end{enumerate}

\section{Related Work and Novelty Boundary}
ORBIT is not a claim to invent optimization of QK/OV products or low-rank quotient geometry.
Knight's recent Riemannian optimizer explicitly treats QK/VO products as rank-constrained
optimization objects while preserving factorized implementations~\cite{knight2026}.
The present work instead targets the setting where RoPE prevents the query--key interaction
from collapsing to one position-independent product. The final paper must update this
section from the branch's dated prior-art gate immediately before submission.

\section{Method}
For a single attention head and RoPE frequency pair $f$, write the relative-position score
contribution as
\begin{equation}
  s_{ij,f}=q_{i,f}^{\top}R_f(\Delta)k_{j,f},\qquad \Delta=i-j.
\end{equation}
Let $C_{Q,f}$ and $C_{K,f}$ be exponential-moving-average $2\times2$ covariances of the
unrotated query and key coordinates. ORBIT defines
\begin{align}
M_{Q,f} &= \mathbb{E}_{\Delta}\left[R_f(\Delta)C_{K,f}R_f(\Delta)^\top\right],\\
M_{K,f} &= \mathbb{E}_{\Delta}\left[R_f(\Delta)^\top C_{Q,f}R_f(\Delta)\right].
\end{align}
The expectation uses a frozen logarithmic displacement grid. A Muon candidate direction
$U_Q,U_K$ is computed first. For exponent $p$,
\begin{align}
\widetilde U_{Q,f} &= M_{Q,f}^{-p/2} U_{Q,f},\\
\widetilde U_{K,f} &= M_{K,f}^{-p/2} U_{K,f}.
\end{align}
A single joint scalar then restores
\begin{equation}
\|\widetilde U_Q\|_F^2+\|\widetilde U_K\|_F^2
=
\|U_Q\|_F^2+\|U_K\|_F^2,
\end{equation}
preventing an aggregate step-size increase from masquerading as a geometric gain. All
matrix operations unique to ORBIT are batched $2\times2$ eigendecompositions.

\section{Experimental Protocol}
The discovery and confirmatory campaign inherits the ASTRO paper harness's pinned
FineWeb-Edu corpus, common warmup--cosine LR factor, GPT-style weight-decay routing,
environment fingerprinting, and faithful AdaMuon reference. The research model is
GPT-2-sized but replaces learned absolute positional embeddings with RoPE; this is required
by the hypothesis rather than presented as a model improvement. Every optimizer sees the
same model and data.

Hyperparameter discovery uses the same number of free tuning dimensions for ORBIT and Muon.
The ORBIT functional exponent, displacement grid, and condition cap are frozen mechanism
choices. Five held-out seeds are evaluated only after discovery configurations are frozen.
The branch also defines 2700-step and 355M transfer phases.

\section{Results}
\IfFileExists{generated_results.tex}{\input{generated_results.tex}}{No confirmatory results yet.}

\IfFileExists{figures/confirm_training_curves.pdf}{%
\begin{figure}[t]
\centering
\includegraphics[width=0.78\linewidth]{confirm_training_curves.pdf}
\caption{Held-out training trajectories.}
\label{fig:curves}
\end{figure}}{}

\IfFileExists{figures/confirm_efficiency.pdf}{%
\begin{figure}[t]
\centering
\includegraphics[width=0.67\linewidth]{confirm_efficiency.pdf}
\caption{Validation loss versus measured wall-clock time.}
\label{fig:efficiency}
\end{figure}}{}

\section{Mechanism Ablations}
The required controls are \texttt{orbit\_identity}, \texttt{orbit\_norope}, and
\texttt{orbit\_diag}. Identity tests whether the implementation is merely another Muon
recipe; no-RoPE tests whether ordinary Q/K covariance adaptation explains the effect; and
diagonal removes the two-dimensional phase coupling while preserving per-frequency scaling.

\IfFileExists{figures/ablations.pdf}{%
\begin{figure}[t]
\centering
\includegraphics[width=0.67\linewidth]{ablations.pdf}
\caption{ORBIT mechanism ablations.}
\label{fig:ablations}
\end{figure}}{}

\section{Scaling and Production Relevance}
ORBIT changes only training-time optimization. Standard Q/K/V/O weights are retained, so
there is no additional inference state or compute. A successful 124M result is not by itself
sufficient: the branch requires a longer-horizon test and a 355M transfer before a strong
performance claim. A subsequent modern-architecture gate must test GQA and QK-normalization.

\section{Limitations and Falsification Criteria}
The local metric assumes that small query/key perturbations can be usefully characterized by
second-order frequency-pair statistics. It is an approximation to attention-logit function
space, not an exact natural gradient. The method should be rejected as a paper contribution
if gains disappear under held-out seeds, if the no-RoPE or identity controls match ORBIT, or
if the wall-clock overhead is not defensible for production training.

\section{Conclusion}
To be written after the confirmatory campaign. The conclusion must report failures and
negative ablations with the same prominence as positive results.

\appendix
\section{Reproducibility}
All runs are append-only JSONL records containing the optimizer configuration, random seed,
code digest, Python/package environment, GPU model, wall-clock time, validation loss,
training curve, and peak CUDA allocation. Four Colab workers write independent shard files;
merging refuses mixed implementation or environment fingerprints.

\bibliographystyle{plain}
\bibliography{references}
\end{document}
